\subsection{Relation to earlier work}

Reversible clustering processes and their product-form equilibrium
laws are classical; see Kelly~\cite[Chapter~8]{Kelly1979} and
Whittle~\cite{Whittle1986}. Durrett, Granovsky, and
Gueron~\cite{DurrettGranovskyGueron1999} characterize reversibility
for a broad class of coagulation--fragmentation processes and study
the resulting equilibrium distributions. Berestycki and
Pitman~\cite{BerestyckiPitman2007} relate Gibbs laws on partitions
to recursive fragmentation and its coalescent time reversal.
The stationary law used here belongs to this established setting.

The solvable family of Granovsky and
Erlihson~\cite{GranovskyErlihson2009} contains the present model.
In their notation, set
\[
 N=n,\qquad a=0,\qquad b=2/n,\qquad\phi(1,1)=\beta.
\]
Their Proposition~1 gives the autonomous count process, and their
fragmentation formulas give the same unequal- and equal-fragment
rates as~\eqref{eq:fragmentation-rates}. Their Section~2.3 announces
equality of the full-process and count gaps and gives the bounds
\[
 \beta+2/n\le\gamma_Q\le\beta+2-2/n.
\]
The displayed transition formula used there concerns initial laws
with Gibbs conditional distributions at fixed count. Our proof of
the gap equality instead controls every orthogonal sector of every
polymer observable. We use this independent comparison to obtain
the uniform error estimate, the critical limit, and the graph
extensions. The gap equality itself is attributed to that earlier
work.

Hoze and Holcman~\cite{HozeHolcman2018} analyze the same constant-kernel
model after a time change: take their fragmentation parameter to be
$\beta n/2$ and multiply their generator by $2/n$. They obtain
equilibrium cluster statistics and quantities for tagged clusters.
D'Orsogna, Lei, and Chou~\cite{DOrsognaLeiChou2015} study first
assembly and equilibration in finite stochastic clustering systems,
including restrictions on the largest allowed cluster. The
all-observable spectral estimate and the mixing statements here
concern the unrestricted partition-valued chain with the stated
normalization.

Our one-dimensional arguments use positive-vector estimates for
birth--death chains, developed in work including
Zeifman~\cite{Zeifman1991} and Granovsky and
Zeifman~\cite{GranovskyZeifman2000,GranovskyZeifman2005}.
The derivative identity is a finite unweighted instance of the
intertwining relations of Chafa\"i and Joulin~\cite{ChafaiJoulin2013}.
Filmus~\cite{Filmus2016} gives the orthogonal harmonic structure on
Boolean slices. We prove the finite identities used in the argument
so that the cut-to-polymer comparison and its boundary conventions
are explicit.

The cutoff step uses the birth--death hitting-time theorems of Chen
and Saloff-Coste~\cite{CSC2015}. The calculation here supplies the
endpoint means and variances for the moving equilibrium scale
$K\asymp\sqrt n$. More general concentration and cutoff results for
density-dependent chains appear in Barbour, Brightwell, and
Luczak~\cite{BBL2022,BBL2026}. A direct application of a
fixed-density theorem would require checking both the
$n$-dependent rescaled rates and the all-monomer starting state,
whose rescaled density diverges. The passage-time proof handles
these features directly.

Spatial coagulation--fragmentation has a separate literature.
Bernardin and Toninelli~\cite{BernardinToninelli2012} study a
one-dimensional interval process with a dynamical phase transition
and mixing estimates. Ioffe and T\'oth~\cite{IoffeToth2020} study
split-and-merge limits for stationary random stirring on a lattice
torus. Our graph model retains exchangeable
reaction rates and adds exclusion motion. It is related to graph
gap reductions for exclusion with reservoirs, including
Salez~\cite{Salez2023,Salez2024}, and to the one- and two-particle
identities of Kim and Sau~\cite{KimSau2026}. The reservoir models
have fixed sitewise creation and annihilation rates, whereas our
deletion rate is $(|S|-1)/n$ and the reversible law is not a product
law. Here the higher reaction sectors already dominate the count
gap, leaving only the one-particle graph calculation in the exact
formula. The proof therefore does not require the general
interchange-process gap theorem of Caputo, Liggett, and
Richthammer~\cite{CaputoLiggettRichthammer2010}.

There are also established growing-kernel and nonreversible
comparisons. Jiang~\cite{Jiang2015} diagonalizes the Metropolis
transposition chain with Ewens equilibrium. We record its spectral
gap in the present time units as a known baseline for a different
multiplicative coagulation rule. Antisymmetric perturbations that
preserve equilibrium can accelerate convergence; Hwang, Hwang-Ma,
and Sheu~\cite{HwangHwangMaSheu2005} give a spectral treatment for
diffusions. Our four-state circulation verifies strict acceleration
on the actual reaction graph of the linear-polymer chain.

Deterministic coagulation--fragmentation equations require different
arguments. Aizenman and Bak~\cite{AizenmanBak1979} prove convergence
through decay of free energy. Recent work of Bhattacharya and
Throm~\cite{BhattacharyaThrom2026} obtains convergence without
detailed balance for perturbations of constant continuous kernels.
The finite graph and kernel results below concern their explicitly
defined Markov generators. They do not supply a convergence theorem
for general spatial coagulation--fragmentation equations.
